\section{Choose which coefficient information to retain}
The transform has rearranged the block's information. Quantization now discards
some of it. A quantizer maps several possible coefficient values to one integer
level. The decoder receives that level and reconstructs one coefficient value;
it cannot infer which original value produced it. This is where the encoder's
local reconstruction begins to differ deliberately from the source.

Use $C$ for a forward coefficient, $L$ for its transmitted level, $Q$ for the
quantizer parameter and $R$ for its reconstructed coefficient. Keeping these
names separate matters: a level is not an inverse-transform input until it has
been inverse-quantized. Likewise, an INTRA DC code byte is not the coefficient
1024 that the transform expects for a flat block of sample value 128.

For this construction exercise, choose the reproducible forward policy
\[
 m=\min\left(127,\left\lfloor\frac{|C|}{2Q}\right\rfloor\right),
 \qquad L=\begin{cases}-m&C<0,\\m&C\geq0.\end{cases}
\]
Here $1\leq Q\leq31$. The cap keeps ordinary nonzero levels within the baseline
signed range $[-127,127]$. Zero is represented by omission when coefficients
are eventually coded. It is still stored explicitly in this exercise's array.
A coefficient with magnitude below $2Q$ maps to zero. That interval around zero
is the \emph{dead zone}: small changes inside it do not survive the quantizer.
At $Q=8$, both 15 and $-15$ disappear, while 16 and $-16$ become 1 and $-1$.

This forward policy is a teaching choice, not a decoder requirement. A complete
encoder can choose other legal levels to trade picture error against stream size.
The kernel exercise checks the stated policy so that intermediate observations
are reproducible. Later whole-encoder grading must allow other choices when the
stream meets its syntax, reconstruction, quality and inter-picture requirements.

Baseline reconstruction has a fixed meaning. An ordinary zero level reconstructs
zero. For a nonzero level, form a positive magnitude, restore its sign, then clip:
\[
 a=Q(2|L|+1)-e,\qquad
 e=\begin{cases}1&Q\text{ even},\\0&Q\text{ odd},\end{cases}
 \qquad R=\operatorname{clip}_{[-2048,2047]}\!\left(\operatorname{sign}(L)a\right).
\]
Do not evaluate this nonzero expression for $L=0$: at $Q=8$ it would manufacture
seven where there should be zero. Do not apply the even-$Q$ adjustment after
restoring the sign. A negative level must receive the same positive magnitude
as its positive counterpart before the sign is applied.

\begin{center}
\tagpdfsetup{table/header-rows={1}}
\begin{tabular}{rrrrl}
\textbf{$Q$} & \textbf{$C$} & \textbf{$L$} & \textbf{$R$} & \textbf{What happened}\\
8 & 15 & 0 & 0 & Dead-zone removal\\
8 & 16 & 1 & 23 & Even-parameter adjustment\\
8 & $-16$ & $-1$ & $-23$ & Sign restored after magnitude\\
8 & 32 & 2 & 39 & Next retained level\\
7 & 16 & 1 & 21 & Odd parameter, no subtraction\\
1 & 2047 & 127 & 255 & Level cap limits reconstruction
\end{tabular}
\end{center}
The last row exposes a limitation of the chosen forward policy at very small
parameters: a large coefficient is capped even though the original coefficient
fits the transform's twelve-bit domain. Increasing detail retention cannot be
reasoned about from $Q$ alone when a level cap is active. Do not use a successful
kernel test as evidence that every parameter gives acceptable video quality.

\section{Give baseline INTRA DC its own path}
For an INTRA block, coefficient zero describes the block's absolute average.
Use a separate eight-unit reconstruction spacing. In this exercise, select
\[
 d=\operatorname{clip}_{[1,254]}\!\left(\left\lfloor\frac{C+4}{8}\right\rfloor\right),
 \qquad R_{\mathrm{DC}}=8d.
\]
The stored DC code is $d$, except that $d=128$ is stored as 255. Thus code 255
means reconstructed coefficient 1024, not 2040. Neither code zero nor code 128
is emitted. The ordinary $Q$ formula still applies to the other 63 coefficients.
For an INTER block, coefficient zero uses that ordinary formula too.

\begin{center}
\tagpdfsetup{table/header-rows={1}}
\begin{tabular}{rrrr}
\textbf{INTRA $C[0]$} & \textbf{Selected $d$} & \textbf{Stored DC code} & \textbf{$R[0]$}\\
0 & 1 & 1 & 8\\
1016 & 127 & 127 & 1016\\
1020 & 128 & 255 & 1024\\
1024 & 128 & 255 & 1024\\
2047 & 254 & 254 & 2032
\end{tabular}
\end{center}
The first row is intentional. The lowest baseline INTRA DC reconstruction is
eight. It is not an ordinary zero coefficient. An all-zero INTRA input can
therefore retain a nonzero reconstructed DC while having no coded AC levels.
That distinction will determine how block-presence flags are formed later.

\section{Trace the first row before writing vector code}
With $Q=8$ and INTRA enabled, use the supplied input beginning
\[
 (1024,15,16,-16,31,32,-32,2047).
\]
The two records begin
\[
 \begin{aligned}
  L&=(255,0,1,-1,1,2,-2,127,\ldots),\\
  R&=(1024,0,23,-23,23,39,-39,2039,\ldots).
 \end{aligned}
\]
The next coefficient, $-2048$, becomes
level $-127$ and reconstructed coefficient $-2039$. There are seven nonzero AC
levels in the supplied complete block. The DC code is excluded from that count.

Write down both records. If you inspect the first array as coefficients, the
flat block appears to have DC 255. If you feed the second array directly to a
bit writer, its values exceed the ordinary level range. The two arrays belong
to different sides of the quantization boundary even though both use halfword
storage in the program.

Inverse quantization does not recover the original coefficients. Its purpose is
to reproduce the coefficient block that a conforming decoder will use. The local
inverse transform must consume this reconstructed block. Using the original
forward block would make the encoder predict from pictures that the decoder
never reconstructed.
